\chapter{Appendix 12: Scenario Mu - Regenerative Extraction}

\section{The Legacy Paradigm \& Its Failures}
Organic biomass is landfilled to produce methane or burned openly, while human biological waste is flushed into centralized chemical treatment plants using purified drinking water. This breaks the thermodynamic loop, forcing reliance on imported petrochemical fertilizers and squandering localized nitrogen and carbon sinks.

\section{The Incremental Automation Pathway}
\textbf{Phase 1: Manual Trust.} Stewards visually inspect compost temperatures and manually schedule pyrolysis burns based on weather reports.

\textbf{Phase 2: AI-Assisted Policy.} Thermocouples provide alerts when kilns deviate from optimal curves or when compost fails to reach thermophilic temperatures.

\textbf{Phase 3: Autonomous Cryptographic Enforcement.} Layer 5 cryptographically locks compost from agricultural use unless L2 telemetry provides an unbroken cryptographic proof of pathogen destruction ($>55^\circ\text{C}$ for 3 days), and automatically suspends burn intents based on external weather oracles.

\section{The Human Boundary (Limits of Automation)}
While automation guarantees the thermal curve, the physical handling of biological waste and the precise loading of the retort kiln remain deeply physical, high-risk tasks requiring trained human stewards to ensure proper inoculation and safety protocols are followed.

\section{Orchestration Mechanics (The "How")}
The orchestrator tracks continuous L2 heat telemetry. If the thermophilic pathogen lockout policy is satisfied, the data is hashed and anchored to the Ledger (Chapter 3) to mint Value Tokens for carbon sequestration. The system also requires ingestion of external municipal weather APIs to enforce burn bans (Triggering Rule 4: Decentralized Oracles for external data without introducing centralized failure points are missing from the core architecture).
